\section{系统与方法}

\subsection{端到端执行顺序}

该系统的执行顺序决定了各分值的含义。图~\ref{fig:pipeline} 从配置解析开始：系统在 seed 为空时生成随机数，再由 Transformers 同时设置 Python、NumPy 与 PyTorch 随机状态。\code{Model} 随后从同一模型标识和可选 \code{model\_commit} 加载 tokenizer、可选 processor 与模型，发现可消融模块，并初始化零增量 LoRA。模型和文本处理器没有彼此独立的 revision 参数。

\begin{figure*}[t]
  \centering
  \resizebox{\textwidth}{!}{%
  \begin{tikzpicture}[
    node distance=8mm and 8mm,
    box/.style={draw, rounded corners=2pt, align=center, minimum height=11mm,
      minimum width=29mm, fill=blue!4},
    data/.style={box, fill=green!6},
    score/.style={box, fill=orange!8},
    artifact/.style={box, fill=purple!6},
    flow/.style={-{Latex[length=2.2mm]}, thick},
    state/.style={-{Latex[length=2.2mm]}, thick, dashed},
    optional/.style={-{Latex[length=2.2mm]}, thick, dashed}
  ]
    \node[box] (settings) {解析 \code{Settings}\\固定或生成随机种子};
    \node[box, right=of settings] (model) {同一模型 revision\\加载处理器与模型\\发现模块并置零 LoRA B};
    \node[data, right=of model] (directiondata) {加载 good/bad 数据\\探测 batch 与回答前缀};
    \node[score, right=of directiondata] (baseline) {加载 scorer 评估数据\\在未消融模型上\\建立 baseline};
    \node[box, right=of baseline] (residuals) {方向数据的首次生成步\\提示末位残差};
    \node[box, right=of residuals] (directions) {均值差方向\\可选 projected 处理};

    \node[box, below=15mm of directions] (trial) {Optuna trial\\方向范围、方向层与\\各组件层权重核};
    \node[score, below=15mm of baseline] (scoring) {重置后施加消融\\由 scorer 计算\\当前 trial 分值};
    \node[score, left=of scoring] (pareto) {multivariate TPE\\多目标 Pareto 集\\用户选择 trial};
    \node[artifact, left=of pareto] (export) {保存 LoRA adapter\\或合并模型};
    \node[artifact, left=of export] (reproduce) {可选复现包\\config、requirements、JSON、\\journal 与哈希};

    \draw[flow] (settings) -- (model);
    \draw[flow] (model) -- (directiondata);
    \draw[flow] (directiondata) -- (baseline);
    \draw[flow] (baseline) -- (residuals);
    \draw[flow] (residuals) -- (directions);
    \draw[flow] (directions) -- (trial);
    \draw[flow] (trial) -- node[above, sloped, font=\scriptsize]{评估} (scoring);
    \draw[flow] (scoring) -- (pareto);
    \draw[flow] (pareto) -- (export);
    \draw[optional] (export) -- node[above, font=\scriptsize]{仅 Hub 上传且满足条件} (reproduce);
    \draw[state] (baseline) -- node[right, font=\scriptsize]
      {缓存基线供比较} (scoring);
  \end{tikzpicture}%
  }
  \caption{自动化系统在固定源码基线中的端到端执行顺序。Evaluator 先加载 scorer
  自有评估数据并缓存未消融 baseline，方向计算完成后各 trial 才调用 scorer 评分；
  方向数据与评估数据角色分离。本地保存不会自动生成复现包，虚线表示状态依赖或
  满足数据固定、内置可复现 scorer 与 Hub 上传等条件后的可选分支。}
  \label{fig:pipeline}
\end{figure*}


方向发现数据在 scorer 之前加载。系统先取得 good/bad 提示，探测可行 batch size，并在用户没有给出 \code{response\_prefix} 时比较生成回答的共同前缀。此后 \code{Evaluator} 才加载插件及各插件自己的评估提示，并在尚未消融的模型上缓存 baseline。默认配置中，方向使用训练切片，KeywordRate 与 KLDivergence 使用测试切片；两类数据承担不同角色，不能统称为“训练集”或“评估集”。表~\ref{tab:components} 给出每个步骤的源码定位。

\begin{table*}[t]
  \caption{固定提交中的实现组件、职责与源码定位}
  \label{tab:components}
  \centering
  \scriptsize
  \begin{tabularx}{\textwidth}{@{}p{18mm}p{42mm}Yp{47mm}@{}}
    \toprule
    组件 & 主要职责 & 本文使用的 claim & 固定提交中的定位 \\
    \midrule
    \code{main.py} & 编排配置、数据、残差、搜索、trial 选择与导出 & 执行顺序、TPE、Pareto 与复现包触发条件 & 行 267--270、410--578、591--735、772--850、1113--1172、1239--1258 \\
    \code{model.py} & 加载模型、发现层模块、构造 LoRA、生成响应和残差、施加方向消融 & 混合层映射、首 token 残差、方向插值、层权重核、FULL 低秩近似 & 行 66--80、189--234、381--449、461--619、630--785 \\
    \code{evaluator.py} & 加载 scorer 插件、建立 baseline、保持多目标顺序 & 方向数据与评估数据分离；目标名、值、方向同序 & 行 34--45、47--109、178--266 \\
    \code{plugin.py} & 动态加载内置或外部插件并提供缓存上下文 & scorer 拥有各自数据和模型服务，外部插件影响复现条件 & 行 46--147、150--184 \\
    \code{scorer.py} 与内置 scorers & 定义分值协议；实现关键词率与首 token KL & 空回答规则、字符串规范化、$D_{\mathrm{KL}}(p_0\|p_\theta)$ 与零 baseline & \code{keyword\_rate.py} 行 10--138；\code{kl\_divergence.py} 行 13--74 \\
    \code{config.py} & 定义 Pydantic 配置、数据集和 scorer 契约 & \code{optimization}、思考前缀、行归一化与搜索参数的归属 & 行 95--125、259--299、337--410 \\
    \code{utils.py} & 加载文本提示、序列化配置与构造上传复现工件 & \code{Prompt} 仅含 system/user；reproduce.json v3 不含 trial 编号 & 行 145--255、557--589、631--735 \\
    \code{reproduce.py} & 消费复现描述并核对环境差异 & \code{--reproduce} 路径用于读取和验证已有复现材料 & 行 35--67、84--213 \\
    \bottomrule
  \end{tabularx}
\end{table*}


\subsection{回答起点的逐层方向}

设层索引为 $l$，有害提示集为 $\mathcal{B}$，无害提示集为 $\mathcal{G}$。$h_l(x)$ 表示模型开始生成回答时，提示末位用于预测首个新词元的残差向量。实现调用只生成一个词元的 \code{generate}，读取首次生成步的 hidden states，再抽取各层张量的最后一个提示位置。由此定义

\begin{align}
\mu_l^{B} &= \frac{1}{|\mathcal{B}|}\sum_{x\in\mathcal{B}} h_l(x), &
\mu_l^{G} &= \frac{1}{|\mathcal{G}|}\sum_{x\in\mathcal{G}} h_l(x), \\
r_l &= \operatorname{norm}\!\left(\mu_l^{B}-\mu_l^{G}\right),
\label{eq:direction}
\end{align}

其中 \(\operatorname{norm}(z)=z/\lVert z\rVert_2\)。普通路径逐 batch 把 32 位浮点（FP32）残差在中央处理器上以 64 位浮点（FP64）累加，求均值后转回 FP32。开启残差几何打印或绘图时，系统则物化 FP32 全量残差并直接求均值。两条路径不能合并描述为共同的 FP64 计算。若启用 winsorization，系统对每个“提示$\times$层”向量沿隐藏分量维计算绝对值分位点并对称截断；它没有跨样本删除提示。

投影式方向先从 $r_l$ 中删去其在无害均值单位方向上的分量。令 $g_l=\operatorname{norm}(\mu_l^G)$，则

\begin{equation}
\widetilde r_l=\operatorname{norm}\!\left(r_l-(r_l^{\top}g_l)g_l\right).
\label{eq:projected-direction}
\end{equation}

为统一表示，启用投影式方向时令 $q_l=\widetilde r_l$，否则令 $q_l=r_l$。trial 可选择逐层使用 $q_l$，也可采样浮点全局索引 $j$。令 $k=\lfloor j\rfloor$、$\alpha=j-k$，全局单位方向为

\begin{equation}
v(j)=\operatorname{norm}\!\left((1-\alpha)q_k+\alpha q_{k+1}\right).
\label{eq:direction-interp}
\end{equation}

因此，浮点索引既不是层号四舍五入，也不是权重矩阵插值。源码中残差数组还包含 embedding 状态，实际取数组时相对 Transformer 层索引偏移一位；该存储细节不改变式~\eqref{eq:direction-interp} 的层间插值含义。

\subsection{组件发现与层权重核}

模型扫描把标准注意力的 \code{self\_attn.o\_proj} 与线性注意力的 \code{linear\_attn.out\_proj} 统一标记为 \code{attn.o\_proj}，并把稠密前馈网络的 \code{mlp.down\_proj} 标记为另一组件。LoRA targeting 仍保存每个真实叶模块全名，因而混合层可以共享搜索参数而不要求物理模块同名。若某层找不到任何受支持投影，系统以断言失败，而不是静默跳过整个架构。

每个组件 $c$ 独立采样最大权重 $M_c$、连续峰值位置 $p_c$、边界权重 $m_c$ 与支撑半径 $d_c$。整数层 $l$ 的权重为

\begin{equation}
\lambda_{c,l}=\begin{cases}
M_c+\dfrac{|l-p_c|}{d_c}(m_c-M_c), & |l-p_c|\le d_c,\\[3pt]
0, & |l-p_c|>d_c.
\end{cases}
\label{eq:kernel}
\end{equation}

图~\ref{fig:weight-kernel} 展示该核。当 $m_c>0$ 时，支撑边界内的值与边界外零值之间存在跳变。$p_c$ 是连续采样值，只有它恰落在整数层上时，某个实际层才严格取得 $M_c$；一般情况下，最近整数层的权重略低于连续核峰值。

\begin{figure*}[t]
  \centering
  \resizebox{\textwidth}{!}{%
  \begin{tikzpicture}[
    x=0.8cm,y=1.0cm,
    block/.style={draw, rounded corners=2pt, align=center, minimum height=10mm,
      text width=32mm, fill=blue!4},
    flow/.style={-{Latex[length=2.2mm]}, thick}
  ]
    \begin{scope}[shift={(4.6,0)}]
      \draw[->] (0,0) -- (7.1,0) node[right] {层 $l$};
      \draw[->] (0,0) -- (0,2.5) node[above] {$\lambda_{c,l}$};
      \draw[dashed, gray] (1.3,0) -- (1.3,0.55);
      \draw[dashed, gray] (3.3,0) -- (3.3,2.05);
      \draw[dashed, gray] (5.3,0) -- (5.3,0.55);
      \draw[very thick, blue, dashed] (0,0) -- (1.3,0);
      \draw[very thick, blue, dashed] (1.3,0.5) -- (3.3,2.0) -- (5.3,0.5);
      \draw[very thick, blue, dashed] (5.3,0) -- (6.7,0);
      \draw[blue, fill=white, very thick] (3.3,2.0) circle (1.8pt);
      \filldraw[blue] (1.3,0.5) circle (1.6pt);
      \filldraw[blue] (5.3,0.5) circle (1.6pt);
      \draw[blue, fill=white, very thick] (1.3,0) circle (1.6pt);
      \draw[blue, fill=white, very thick] (5.3,0) circle (1.6pt);
      \foreach \x/\y in {1/0,2/1.025,3/1.775,4/1.475,5/0.725,6/0}
        \filldraw[black] (\x,\y) circle (1.5pt);
      \node[below] at (1.3,0) {$p_c-d_c$};
      \node[below] at (3.3,0) {$p_c$};
      \node[below] at (5.3,0) {$p_c+d_c$};
      \node[left] at (0,2.0) {$M_c$};
      \node[left] at (0,0.5) {$m_c$};
      \node[align=center, font=\scriptsize] at (3.4,2.65)
        {虚线为连续核；黑点为整数层实际采样\\
        $p_c\notin\mathbb{Z}$ 时实际层不取得 $M_c$};
      \node[align=left, font=\scriptsize] at (7.7,1.1)
        {$m_c>0$ 时边界处跳变\\各组件独立采样一组参数};
    \end{scope}

    \begin{scope}[shift={(0,-4.3)}]
      \node[block] (directions) {相邻层有效方向\\
        $q_{\lfloor j\rfloor},q_{\lceil j\rceil}$};
      \node[block, right=8mm of directions] (interp) {线性插值后归一化\\
        $v=\operatorname{norm}($\\$\operatorname{lerp}(\cdot))$};
      \node[block, right=8mm of interp] (project) {秩一投影项\\
        $v(v^{\top}W_{c,l})$};
      \node[block, right=8mm of project] (updated) {组件与层特定更新\\
        $W'_{c,l}=W_{c,l}-$\\$\lambda_{c,l}v(v^{\top}W_{c,l})$};
      \draw[flow] (directions) -- node[above, font=\scriptsize]
        {浮点 direction index} (interp);
      \draw[flow] (interp) -- (project);
      \draw[flow] (project) -- node[above, font=\scriptsize]
        {$\lambda_{c,l}$} (updated);
      \node[draw, dashed, fit=(project)(updated), inner sep=4mm,
        label={[font=\scriptsize]below:FULL 模式随后恢复行范数并以随机低秩分解近似差值}] {};
    \end{scope}
  \end{tikzpicture}%
  }
  \caption{上：组件 $c$ 的连续层权重核及其整数层采样；不同组件独立采样核参数。
  下：浮点方向索引、矩阵投影与层权重的关系。FULL 行归一化在秩一更新之后
  引入非线性，因此最终 LoRA 只近似完整差值。}
  \label{fig:weight-kernel}
\end{figure*}


\subsection{矩阵投影与 LoRA 近似}

对无行归一化的基础路径，令 $W_{c,l}\in\mathbb{R}^{d_{\mathrm{out}}\times d_{\mathrm{in}}}$，选定单位方向 $v\in\mathbb{R}^{d_{\mathrm{out}}}$，则更新为

\begin{equation}
W'_{c,l}=W_{c,l}-\lambda_{c,l}v(v^{\top}W_{c,l}).
\label{eq:weight-edit}
\end{equation}

系统用 $B=-\lambda_{c,l}v$ 与 $A=v^{\top}W_{c,l}$ 表示该秩一增量。若基础权重为 4-bit，代码先把它反量化为 FP32 再计算。每个 trial 开始前，系统清零所有 LoRA B 矩阵，因此新的参数不会叠加在前一个 trial 上。

预归一化（pre-normalization, PRE）模式先对 $W$ 的每一行归一化。设行范数对角矩阵为 $D_s$、$\widehat W=D_s^{-1}W$，则代码实际设置

\begin{equation}
A=v^{\top}\widehat W,\qquad B=-\lambda D_s v.
\label{eq:pre}
\end{equation}

这一步把原行范数乘入 B，而不是简单复用式~\eqref{eq:weight-edit} 的两个因子。全量归一化（full normalization, FULL）模式在 $\widehat W$ 上先施加秩一更新，再把调整后矩阵逐行归一化并恢复原始行范数，最后减去原 $W$ 得到完整差值。代码在每个模块分解前重设随机种子，调用 \code{torch.svd\_lowrank}，并截取配置 rank 写入 LoRA。该过程是固定 seed 的随机低秩近似；有限 rank 的 adapter 不保证严格复现完整差值，也不保证合成后的每行范数完全等于原矩阵。
